% TeX root = ../Main.tex

% First argument to \section is the title that will go in the table of contents. Second argument is the title that will be printed on the page.
\section[Question -- {\it Detecting keyframes in Video}]{}

\subsection{Detecting keyframes in Video}
The all approach can be devided in 3 main features, to detect all the different situation in a video sequence.
\begin{itemize}
    \item Macro (global changes)
    \item Micro changes
    \item Optical flow magnitude for movements 
\end{itemize}

The all algorithm is completed by using a minimum gap constraint to avoid considering a number k of successive frames, 
to not face the redundancy problem.

\subsection{Macro}
To have a general and comprehensive understanding of the scene, the color histogram in HSV can be analized. This approach
is not complex and allows to detect scene interrupts and macro changes in the video sequence. The two histograms are compared
with Chi square distance.

\subsection{Micro}
It may happen that the color histogram in two consequent scenes is similar, but the objects displacement is completely different.
This issue may be overcame by the use of an edge ratio system: a canny edge detector algorhitm is used to find all edges in a scene
and the amount of edges in two scenes is compared finding a change ratio.

\subsection{Optical Flow}
Finally, to detect the beginning of actions, the optical flow magnitude can be used $M = \frac{1}{N}\sum \sqrt{u^2 + v^2}$. 
LK comparing algorithm is used to compare consecutive frames and it can detect in two scenes some movements (mouth moving 
while talking), where the scenes appear similar to the Macro analyzer.


